\section{A language whose runtime functions are qubits}
The underlying operation is a one-selector quantum multiplexor in the sense of Shende, Bullock and Markov \cite{shende}. QAF exposes its selector as a named wire that a program can invoke and later probe.
QAF has an explicit separation between a program that describes an experiment and the quantum systems in that experiment. The \texttt{fn} binding associates a function name with a qubit wire and a statically specified pair of unitary targets $U_0,U_1$. An unnamed wire can also be invoked with the default pair $(I,X)$. The operational action is
\[
 V=\ket0\bra0\otimes U_0+\ket1\bra1\otimes U_1 .
\]
Since $V^\dagger V=I$, this is a valid transformation on a superposed, mixed or entangled control. The implementation acts on the joint register. Naming a wire twice creates aliases of the same system. Using it as both selector and target in one invocation fails the distinct-wires check. The two operation descriptions belong to the invocation interface. They are classical code and additional information beyond the physical qubit. Relative phase matters: replacing $U_1$ by $-U_1$ changes the coherent control even though the isolated target gates give the same channel. The no-programming theorem requires orthogonal program states for distinguishable exact unitary channels in a deterministic universal processor \cite{noprogram}. Two program labels can select two supplied alternatives; an arbitrary unknown unitary cannot be retrieved from the amplitudes of a single qubit.

The working grammar is reproduced in Appendix~\ref{sec:grammar}. TatSu compiles the grammar to a parsing-expression-grammar parser \cite{tatsu}. The language follows the design guidance of the Zen of Python \cite{zen} through visible distinctions between state preparation, physical application, measurement, classical circuit construction and simulation-only inspection. Numerical expressions, class methods, recursion, mutable local bindings and finite loops describe amplitudes, gate entries or circuits. A \texttt{def} is a classical routine and a \texttt{proc} builds an instruction sequence. The runtime quantum-function type is the wire type. Extending the identification to every classical routine, including a routine that returns an arbitrary integer or fails to terminate, would change the meaning of either programming functions or qubits.

\inputminted[fontsize=\footnotesize]{qaf}{../../examples/prefix/recursion.qaf}
The recursive \texttt{sign} method runs for each requested basis label in the initial-state kernel. The \texttt{fan} procedure expands at compile time, yielding a Hadamard and two CNOTs in total with the explicit Hadamard in the main program. The final distribution is supported on $000$ and $111$ with probability $1/2$ each. A depth and step guard admit bounded executions and report exhaustion. In particular, a divergent recursive function is a meaningful program whose evaluation fails; it cannot denote an additional physical qubit. The grammar also supports unitary kernels, Kraus families, sequential and tensor composition, unitary adjoints, coherent selection, and before/after wrappers for slots.

\inputminted{qaf}{../../examples/prefix/functions.qaf}
The three definitions above implement three distinct constructions. The diagonal gate applies the spectral function $\exp(-i\pi Z/3)$. The two Kraus operators dephase in the $Z$ basis. The wrapper conjugates the channel with Hadamards, producing dephasing in the $X$ basis. The initial $\ket+$ is fixed by that channel, and the final Hadamard gives computational zero with probability $\cos^2(\pi/3)=1/4$. This program exercises a scalar function of an operator, a channel acting on a state and a higher-order transformation acting on a channel.

\subsection{Native execution and its reference semantics}
For a target list $w=(w_0,\ldots,w_{k-1})$ let $b[w\leftarrow x]$ replace its bits in the given order, with wire zero the leftmost displayed bit. A gate kernel produces
\[
 f'(b)=\sum_{x\in\B^k}U_{b_w,x}\,f\bigl(b[w\leftarrow x]\bigr).
\]
This is equation~\eqref{eq:isometry} applied to the usual tensor embedding of the gate. A CUDA thread computes one output amplitude, using local wire masks and its selected small matrix. The new state uses separate device storage so simultaneous updates cannot read partially written amplitudes. Scalar class fields and arithmetic in state and operation definitions are emitted to CUDA device functions. NVRTC compiles those definitions through CuPy, including supported loops and recursion \cite{cupy}. Parsing, staging, allocations and kernel launches run on the host. Query outputs cross the device boundary when requested.

An instrument introduces joint classical--quantum sectors $\rho_c$ with $\sum_c\Tr\rho_c=1$. Applying Kraus operators performs $\rho_c\mapsto\sum_a K_a\rho_c K_a^\dagger$. Measuring wire $q$ and overwriting bit $j$ maps all incoming sectors to the corresponding outgoing labels through the two projectors $P_b^{(q)}$. When different old labels lead to the same new one their masses add. A conditional selects input sectors at entry, executes its body on the selected part, and recombines it with the untouched part. This remains valid when the body overwrites its own condition bit. The CPU reference instead tracks normalized branch wavefunctions with weights. Expanding each sector as $\rho_c=\sum_{i:c_i=c}p_i\ket{\psi_i}\bra{\psi_i}$ proves by induction on instructions that the two ideal rules agree. On a physical device a single invocation samples an outcome. Exact branch probabilities are a simulator service.

For $n$ quantum wires, $m$ classical bits and local width $k$, pure-state storage per complex double buffer is $16\,2^n$ bytes, with $O(2^{n+k})$ work per dense gate. The general sector representation uses $16\,2^m4^n$ bytes per buffer. Applying $r$ dense Kraus operators uses $O(2^m4^n r4^k)$ scalar work with the direct kernel. Gate generation, query materialization and temporary buffers can raise peak memory. An NVIDIA RTX A5000 Laptop GPU executed both paths. The portable CPU path shares the same mathematical convention. Norm and Gram tests use absolute tolerance $10^{-10}$. Double precision and checks at a tolerance implement numerical predictions, with no general proof of correctly rounded evaluation for every program.

\subsection{A lossless Lean embedding}
The translator obtains a concrete TatSu derivation with labeled productions, terminals and source trivia. A Lean \texttt{Source} tree represents that derivation. A separate inductive \texttt{Term} gives a programmable embedding that users can build with Lean definitions before evaluation. The maps exchange the respective token, end, argument and rule constructors. The central proofs are short enough to reproduce directly from the checked file \cite{lean}.
\inputminted[firstline=34,lastline=38,fontsize=\footnotesize]{lean4}{../../src/qaf/lean/QAF.lean}
Structural induction proves both inverse identities on all values of the respective inductives. A spelling witness lifts the first identity to original UTF-8 source text. Preservation of any meaning assigned to the source tree follows by substitution along equality. The command-line importer evaluates the closed Lean program, decodes it and reparses its spelling to enforce membership in the canonical image of the grammar. The exporter reads back the generated Lean program and compares its tree. Edits, comments, blank lines and line endings thus have a defined round-trip path for accepted inputs. The theorem covers the embedded \texttt{Term} language. Arbitrary Lean libraries with effects or unrelated types have no implied QAF translation. The handwritten interpreter, host parser and CUDA code generator remain outside the kernel-checked proof.

\subsection{Portable execution and backend validation}
QAF's portable circuit subset starts from zeros and includes $H,X,Y,Z,S,T$, CNOT, exact Clifford decompositions of CZ and SWAP, sequences, and selected coherent or classical controls with registered rules. It accepts marginal computational probabilities and classical records. Measurement can be deferred coherently using fresh record and history wires. To overwrite a classical record, move the old value to a fresh history label and compute the new outcome into the cleared label. For nested conditions capture the branch decision in a fresh flag at entry. Orthogonal history states kill cross terms on marginalization. Induction gives the same requested statistics as immediate measurement. This is an isometric dilation and equality of the observable program semantics. The added ancillas, erased source comments and many equivalent gate schedules prevent a general source-text isomorphism to a provider program.

The Quil emitter has a separate reader. Certification checks each logical wire's ordered gate history and the terminal readout map. Identical histories permit reordering independent disjoint gates. In a common total order, slide a first instruction past disjoint instructions until it aligns with its counterpart; induction on the remaining word proves the ideal unitaries equal. Reversing CNOT's wires or swapping classical destinations is caught. Amazon Braket and Cirq adapters inspect their SDK circuit objects, while the IonQ path inspects its intermediate representation. For a compatible Google device an exact interaction-graph embedding can replace CNOT with $H$--CZ--$H$ and validate the circuit against the device. QPU calibration, compiler optimization after handoff and sampling noise are separate questions \cite{braket,cirq,azure}. Exact support for an arbitrary provider instruction must be established before admitting it.

The annealing profile describes binary or spin variables by $E(x)=c+\sum_i h_i x_i+\sum_{i<j}J_{ij}x_ix_j$. The local simulator can enumerate energies on CUDA or CPU. The D-Wave Ocean adapter reads back the binary quadratic model, including offset and variable convention, before sampling \cite{dwave}. Repeated coefficients are combined; a self-pair becomes a linear term for binary variables or a constant for spins. The one-hot example uses $(x_0+x_1+x_2-1)^2$ and has precisely three minimizers of energy zero. Arbitrary gate circuits have no transparent compilation to that annealing profile. A physical embedding and finite anneal need not find a true minimum on each read.

\begin{table}[ht]
\centering\small
\begin{tabularx}{\linewidth}{@{}lXX@{}}
\toprule Platform & Artifact exercised & Extent of evidence\\\midrule
NVIDIA CUDA & Native wavefunction, mixed-register and classical-sector kernels & Local device execution on an RTX A5000 Laptop GPU\\
CPU & Branch reference and exact rational series & Deterministic local execution\\
Amazon Braket & SDK circuit and local simulator & Circuit readback and sampled Bell support\\
Google Cirq & SDK circuit and local simulator & Circuit readback and sampled Bell support\\
Azure Quantum & Quil on hosted Rigetti QVM & 100 shots, 49 of 00 and 51 of 11\\
Azure IonQ & JSON gate payload & Offline construction and readback\\
D-Wave Ocean & Binary quadratic models & Local exact sampler and coefficient comparison\\
\bottomrule
\end{tabularx}
\caption{Execution evidence. The Azure QVM was the remotely executed target. A physical-QPU run would add evidence about a particular device and calibration.}\label{tab:platforms}
\end{table}

\subsection{Experiments and independently checkable predictions}
The following Python program checks populations, local mixing and joint coherence using the reference evaluator. The last assertion distinguishes a pure Bell state from a classical source that emits $00$ or $11$ with equal probability.
\inputminted[fontsize=\footnotesize]{python}{../../examples/check_callable.py}
For a stronger Bell test, reverse the entangler on the output or measure a joint $X\otimes X$ observable. A Bell state gives expectation one while the equally weighted incoherent mixture of $00$ and $11$ gives zero. The hosted 100-shot Bell run measured in the computational basis, so it checked readout correspondence and the expected support without measuring entanglement. No physical QPU benchmark or speedup estimate follows from these counts.

Other deterministic checks exercise teleportation of a state with a complex relative phase in all four outcome branches, Grover search, order finding, the Toffoli decomposition on all eight input columns including phases, and seeded arbitrary-input gates on permuted wires. Reconstruction tests compare Bloch, SIC, Wigner and Choi forms. A process example checks memory under two interventions. The singular-value example compares the dilated circuit block with its closed polynomial. For the Dirichlet box $(0,1)$ in units $\hbar=m=1$, $N=2^n$ interior points give ground energy $2(N+1)^2\sin^2(\pi/[2(N+1)])$. It converges from below to $\pi^2/2$. The associated infinite series example checks exact overlaps against normalized finite QAF preparations. All tests and numerical outputs are supplied with the paper.
